\subsection{The exact comparison for Kodaira-zero fibers}
\label{rankone:section}

For fibers of logarithmic Kodaira dimension zero, the relative
pluricanonical forms determine a single Hodge line on the base.  We
construct that line, compute its boundary orders, and use them to
identify the adjoint section spaces.  Coefficient-one boundary
components introduce mixed Hodge structures, but the line itself
lies in one pure weight grade.

A pure variation of Hodge structures carries a local system and a
varying Hodge filtration. The \emph{highest step} of a summand is its
highest nonzero filtration piece. The \emph{parabolic extension} of
such a line is obtained by local power substitutions making boundary
monodromy unipotent, extension of the unipotent Hodge filtration, and
descent with the resulting rational weights. We will compute those
weights using the actual relative pluricanonical form.

\begin{proposition}[The rank-one comparison]\label{rankone:comparison}
Let $g:(X,\Delta)\to W$ be a fibration with connected fibers from a
smooth compact K\"ahler manifold to a smooth compact manifold in
class $\mathcal C$.  Suppose that $\Delta$ is an SNC boundary with
rational coefficients in $[0,1]$ and that
$\kappa(G,K_G+\Delta_G)=0$ on a very general smooth fiber $G$.
After the modifications of
\Cref{setup:log-modification,setup:flattening}, fix a smooth compact
reference pair $(X_0,\Delta_0)$ and a modification
\[
 \pi:X\longrightarrow X_0,\qquad
 \Delta=\pi_*^{-1}\Delta_0+\operatorname{Exc}(\pi)_{\mathrm{red}},
\]
such that every prime divisor mapped by $g$ into codimension at least two
is $\pi$-exceptional.  Then there are a rational line bundle $M$
and a rational divisor $T$ on $W$ with the following properties.
\begin{enumerate}[label=\textup{(\roman*)}]
\item A positive integral multiple of $M$ is the parabolic highest Hodge
line of a complex direct summand of a pure, real-polarizable variation
with an integral lattice on a dense open subset of $W$.  Its parabolic extension is the line to which the adjoint-positivity
theorem, \Cref{period:adjoint}, will apply.
\item The divisor $T$ is supported on an SNC divisor and satisfies
\begin{equation}\label{rankone:boundary}
 B(g,\Delta)\leq T\leq 1.
\end{equation}
In particular, $T$ is effective.  If every prime over a prime $D$ of $W$
has $\Delta$-coefficient one, then $T_D=1$.
\item In sufficiently divisible degrees $q$, multiplication by the
relative generating forms gives linear isomorphisms
\begin{equation}\label{rankone:sections}
 H^0\bigl(W,q(K_W+T+M)\bigr)
 \simeq H^0\bigl(X,q(K_X+\Delta)\bigr)
 \simeq H^0\bigl(X_0,q(K_{X_0}+\Delta_0)\bigr).
\end{equation}
In a common divisible grading they preserve products and ratios of
sections, and hence the dimensions
of the images of their linear systems.
\item The same comparison holds relatively.  If $h:W\to Y$ and
$f_0:X_0\to Y$ are holomorphic maps with $f_0\pi=hg$, then
\begin{equation}\label{rankone:relative-sections}
 h_*\mathcal O_W\bigl(q(K_W+T+M)\bigr)
 \simeq (f_0)_*\mathcal O_{X_0}\bigl(q(K_{X_0}+\Delta_0)\bigr)
\end{equation}
for sufficiently divisible $q$.  The comparison of ratios also applies
to systems over $Y$ and to their restrictions wherever generic base
change holds. The relative identity remains valid after tensoring by a
line bundle on $Y$, or a rational line bundle in degrees clearing its
denominator.
\end{enumerate}
\end{proposition}

The Hodge construction and its divisorial calculation are independent
of the reference pair.  The condition on $\pi$ is used in the final
section comparison, to descend forms across exceptional centers.

\subsubsection{The root form and its eigenspace}

The cyclic-root construction and its distinguished eigensheaf occur in
\cite[Remark~2.6]{FujinoMori},
\cite[Lemma~5.2]{AmbroBoundary}, and
\cite[Remark~3.9]{FujinoGongyo}. Here the order calculations also
include coefficient-one boundary and keep the character line distinct
from the entire top Hodge space.

Choose a positive integer $m$ clearing the boundary denominators such
that the generic rank of
\[
 g_*\mathcal O_X\bigl(m(K_{X/W}+\Delta)\bigr)
\]
is one.  The same holds in every positive multiple of this degree:
powers of a nonzero section exist, whereas two independent sections
would contradict Kodaira dimension zero.  Its reflexive hull
$\mathcal E_m$ is a line bundle since $W$ is smooth.  A local frame $s$ of
$\mathcal E_m$ determines a meromorphic relative $m$-canonical form on
$X$ over the corresponding base open set.  Two such frames differ by
an invertible base function.
The direct image records $m$-pluriforms. Positivity will come from the
Hodge extension of their rational root; comparing the two extension
orders will supply $T$.

Resolve the horizontal zeros, poles, and boundary, then choose a dense
open set $W^\circ$ on which $g$ is smooth, the resulting divisors are
relatively SNC, and all their strata are smooth over the base.  By
\Cref{setup:log-modification}, new exceptional components have
coefficient one, and the logarithmic Kodaira dimension of the general
fiber stays zero.  On each base chart take the normalized cyclic cover
of $m$-th roots of $s$ as a relative pluriform and resolve it
functorially.  Write $\tau$ for its tautological relative top form and
$d=\dim X-\dim W$.

At a horizontal prime of boundary coefficient $\delta$, let $l$ be
the order of $s$ divided by $m$.
The inequality $l\geq-\delta\geq-1$ shows that $\tau$ has at most
logarithmic poles.  Indeed, at a covering prime of ramification index
$j$, its order is
\begin{equation}\label{rankone:root-order}
 j(1+l)-1.
\end{equation}
This number is integral.  If $l>-1$ it is nonnegative; if $l=-1$ it
equals $-1$.  On the resolved cover let $H$ be the reduced inverse
image of the horizontal primes with $l=-1$.  It includes the
exceptional components above those primes that are needed in the log
resolution.  Logarithmic pullback preserves the bound, and no
additional horizontal poles occur.  We use the cohomology of the
complement of $H$.

\begin{lemma}\label{rankone:eigenspace}
The tautological form spans the root-character part of
$F^dH^d(\widetilde G\setminus H,\C)$ on a general fiber of the cyclic
cover.  In particular, that part has dimension one.
\end{lemma}

\begin{proof}
For a compact K\"ahler SNC pair, logarithmic Hodge--de Rham degeneration
identifies this top Hodge space with its global logarithmic top
forms; it is also the smooth-fiber case of
\cite[Theorem~1.1]{FujinoFujisawaVMHS}, by duality.  Let $\eta$ be
another form of the tautological character.  Then $\eta/\tau$ is
deck-invariant and descends to a meromorphic function $v$ on $G$.

Consider a prime where the logarithmic section $s$ has no zero.
There $l=-\delta$.  For $l>-1$, the integer in
\eqref{rankone:root-order} lies between $0$ and $j-1$.  A pole of
$v$ downstairs would lower the covering order by at least $j$, which
is impossible.  For $l=-1$, the logarithmic pole allowance is already
exhausted, so the same conclusion holds.  Thus the poles of $v$ can
occur only on the positive zero divisor of $s$ as a section of
$m(K_G+\Delta_G)$.  A sufficiently high power $s^N$ absorbs those
poles.  Both $s^N$ and $v s^N$ are then sections of
$Nm(K_G+\Delta_G)$.  The assumption $\kappa=0$ forces them to be
dependent, and hence forces $v$ to be constant.

The argument also applies to a disconnected cover: we retain the whole
cover with its $\mu_m$-action, whose invariant meromorphic functions
descend to $G$.
\end{proof}

The ranks of the Hodge and weight graded bundles are locally constant
on $W^\circ$.  The conclusion of \Cref{rankone:eigenspace}, proved on
a very general fiber, therefore holds throughout this smooth log
locus.  Notice that neither the argument nor
\eqref{rankone:root-order} requires a boundary coefficient to have
the form $1-1/a$ with $a$ an integer.

\subsubsection{Pure weights and descent of the local systems}

We next place a power of the local Hodge line in a global pure
variation.  This requires compatible polarizations and descent,
because the cyclic covers need not form a global family.

First make the local normalized-cover constructions on charts of the
compact base $W$, using the meromorphic forms supplied by the frames
of $\mathcal E_m$.  They extend across the missing base divisor as
finite covers; subsequent resolutions can be projective and
functorial in the analytic category
\cite[Theorem~1.1]{BierstoneMilmanFunctorial}.
The resulting maps to the K\"ahler source are projective, since the
finite covers are projective. Thus their total spaces are K\"ahler
over relatively compact source neighborhoods
\cite[Definition~4.1 and Lemma~4.4]{FujikiDouady}. Properness allows
such a neighborhood to contain the inverse image of a sufficiently
small disk transverse to a base divisor.
The divisor of the tensor is unchanged under multiplication
by a base unit.  Thus the constructions on overlaps are identified
after taking a local root of that unit, and the resolutions and their
exceptional divisor classes are identified as well.

Fix a K\"ahler class $\xi$ on $X$.  On each resolved-cover chart
$q_i:Z_i\to X|_{U_i}$, choose the relatively ample line bundle $A_i$
as follows.  Functoriality matches the blowup center ideals and their
tautological line bundles under the root-change identifications.
Insert identity stages where a center is empty; the corresponding
tautological bundle is trivial there.  A sufficiently weighted tensor
product gives compatible $A_i$, equivariant under deck transformations,
by descent of the centers, blowups, and tautological bundles.  The classes
\[
 \theta_i=Nq_i^*\xi+c_1(A_i)
\]
are K\"ahler on the smooth compact fibers for $N$ sufficiently large;
their restrictions to the compact boundary strata are K\"ahler too.
At the intermediate finite covering stage one uses the pullback
K\"ahler \emph{class}; the literal pulled-back form can degenerate
at ramification.  The relatively ample class supplies positivity
along the projective resolution.  One may choose a single $N$:
carry out the construction over charts $U_i^+$ containing the closures
of finitely many relatively compact charts $U_i$ covering $W$.
Properness makes the inverse images of these closures compact, so
each admits a bound $N_i$ in this class construction.  Taking
$N\geq\max_iN_i$ works on all charts, since increasing $N$ adds
the semipositive class $q_i^*\xi$.  This is why the construction was
extended over the compact base before shrinking to $W^\circ$.
For the invariant K\"ahler-class construction on an equivariant
resolved cyclic cover, see also \cite[Lemma~8.1]{ChenLimits}.

The classes $\theta_i$ come from total-space cohomology, so their
restrictions to the smooth fibers and strata are flat for the
Gauss--Manin connection.  They agree under the overlap
identifications.  Lefschetz decomposition and cup product consequently
give compatible flat real polarizations on the compact-stratum
cohomology systems.  The weight spectral sequence, with its Tate
twists, gives the pure weight grades of the open-fiber cohomology as
subquotients of these systems.  They are real-polarizable: a real Hodge
subvariation of a polarized pure variation is nondegenerate for the
polarization, and a quotient can be identified with an orthogonal
complement.  These operations are compatible with the overlap maps.
The integral structures come from cohomology modulo torsion and the
saturated intersections with the rational weight filtration.  We do
not assert that the real polarizations are rational.

There is a unique weight $w$ for which
\[
 F^d\Gr^W_w H^d(\widetilde G\setminus H,\C)_\chi\ne0.
\]
Here $\chi$ is the tautological character.  Uniqueness follows from
the rank-one assertion and strictness of the weight filtration.  Let
$V_i$ be the pure integral variation given by this weight grade on
the $i$th chart; the chosen line is $F^d(V_i)_\chi$.

\begin{lemma}\label{rankone:descent}
A positive tensor power of these local Hodge lines is the highest
Hodge line of a complex direct summand of a global real-polarizable
integral pure variation on $W^\circ$.
\end{lemma}

\begin{proof}
Refine an overlap so that a root of its transition unit can be
chosen.  It gives an isomorphism between the cyclic-cover families,
and therefore an isomorphism of their integral cohomology systems.
Another choice differs by a deck transformation.  On a triple
overlap the cocycle defect is also a deck transformation.  All these
maps commute with the distinguished $\mu_m$-action.

Put
\[
 \mathbb U_i=
 \bigl((V_{i,\Z}/\text{torsion})^{\otimes m}\bigr)^{\mu_m},
\]
using the diagonal deck action.  The ambiguity and the cocycle defect
act trivially on $\mathbb U_i$.  Hence the $\mathbb U_i$ descend to
an integral local system $\mathbb U$; saturation only changes the
lattice by an isogeny.  The compatible real polarizations, tensorized
and restricted to invariants, polarize this pure variation.

Over $\C$ the character projector $e_\chi$ is a flat Hodge
endomorphism on each chart.  The projector
$e_\chi^{\otimes m}$ restricts to the invariant tensor space,
because $\chi^m=1$.  It commutes with all transition maps and
therefore defines a global complex direct summand with local fibers
$((V_i)_\chi)^{\otimes m}$.  Its highest step is
$(F^d(V_i)_\chi)^{\otimes m}$, a line.  In local rooted notation
this line is generated by $\tau^{\otimes m}$, which is identified
meromorphically with $s$.  Thus no global root of $\mathcal E_m$ and
no global family of cyclic covers is needed.
\end{proof}

\subsubsection{The two extension orders}

Let $M$ denote one $m$th of the parabolic extension of the line just
constructed.  We use $s^{1/m}$ for its local rational frame; the
notation becomes literal after the powers and finite boundary covers
above. We now compare this parabolic extension with the pole bounds
of the original pluriform.

Fix a prime divisor $D$ on $W$ and use a local ordinary base volume
$\omega$ that does not vanish at its generic point. Work on a source
model on which the divisor of the generator, the boundary, and the
inverse image of $D$ have SNC support over that generic point,
using strict transforms plus reduced exceptional boundary as before.
All primes of this prepared source dominating $D$, including the
exceptional primes introduced by its resolution, enter the following
order tests. Put
\[
 l_E=\frac1m\ord_E(s\wedge g^*\omega^m),\qquad
 a_E=\ord_E(g^*D),\qquad \delta_E=\Delta_E,
\]
and
\[
 \alpha_D=\min_{g(E)=D}\frac{l_E+\delta_E}{a_E},
 \qquad \beta_D=\ord_D^M(s^{1/m}).
\]
The notation $s\wedge g^*\omega^m$ means the relative-times-base
identification in the $m$th tensor power of the canonical bundle.
For $q$ divisible by $m$ and clearing the rational data, and a
meromorphic base function $\varphi$, the total pluriform
$\varphi\,\omega^q s^{q/m}$ has boundary-adjusted order
\[
 a_E\ord_D(\varphi)+q(l_E+\delta_E)
\]
at $E$. It therefore satisfies the source boundary bounds at all
these primes if and only if
\begin{equation}\label{rankone:base-bound}
 \ord_D(\varphi)+q\alpha_D\geq0.
\end{equation}
The Hodge frame contributes $q\beta_D$ to the order on the base,
so the candidate boundary correction is $T_D=\alpha_D-\beta_D$.
Computing $\beta_D$ will reduce the required bounds on $T_D$ to
a comparison of two minima. We first justify computing that Hodge
order by logarithmic orders on a semistable cover, without changing
the chosen pure Hodge line. Global consistency of $T$ and descent
across the untested source primes will follow afterward.

\subsubsection{The extension across a transverse disk}

We use the analytic canonical-extension theorem
\cite[Theorem~1.1]{FujinoFujisawaVMHS}.  It applies to a proper
surjective morphism of an analytic SNC pair $(Z,H)$ to a smooth base,
with $H$ reduced, every stratum K\"ahler and dominant, and all strata
smooth off a normal crossing discriminant.  It supplies a
graded-polarizable real variation on compact-support cohomology,
locally free double Hodge/weight grades on the extension, and the
dual direct-image description by $\omega_{Z/B}(H)$.  Admissibility
is established in \cite[Remark~4.8]{FujinoFujisawaVMHS}.
These conclusions apply to proper K\"ahler families; projectivity
is not required.

Take a transverse disk at a general point of a base divisor. Resolve
the cyclic-cover pair together with the central fiber, leaving its
smooth logarithmic locus over the punctured disk unchanged. Near a
central point the resulting map has the form
\[
 t=\varepsilon x_1^{a_1}\cdots x_r^{a_r},
 \qquad H\subset\{x_{r+1}\cdots x_l=0\},
\]
where $\varepsilon$ is a unit and $H$ is the horizontal boundary. Absorbing it
into a vertical coordinate gives a strict toroidal morphism. Compactness
of the central fiber and properness allow us to shrink the disk until
finitely many such charts cover its inverse image. The associated cone
complex is therefore finite. Apply the combinatorial semistable
subdivision theorem
\cite[Theorem~2.7]{AdiprasitoLiuTemkinSemistable}. Its base is one ray,
so the base lattice alteration is realized by $t=u^N$. Normalize this
base change, keeping all components, and realize the projective source
subdivision analytically
\cite[Lemmas~2.14--2.15, 2.18, and~2.31(4)]
{EnokizonoHashizumeSemistable}. Zero-image faces, which record the
horizontal boundary, are allowed in the combinatorial theorem.

The semistable local form is
\[
 u=y_1\cdots y_q,
 \qquad H\subset\{y_{q+1}\cdots y_k=0\};
\]
see \cite[Lemma~2.21(4)]{EnokizonoHashizumeSemistable}.
Here $H$ includes every horizontal exceptional component in the
transformed boundary. The total space is smooth, the central fiber is
reduced, and its union with $H$ is SNC. On every horizontal
intersection stratum the displayed monomial is nonconstant;
properness makes its image the whole disk. The composite to the
initial K\"ahler total space is projective, so the new total space and
its smooth strata are K\"ahler after shrinking the disk
\cite[Lemma~4.4]{FujikiDouady}.

These toroidal modifications preserve the open pair on the punctured
fibers. Its cohomology is therefore the pullback of the original
open-fiber cohomology, including the full direct sum if the cover
splits. Transport the original $\mu_m$-action through this
identification. The deck projector then acts on the variation and,
by functoriality, on its canonical extension; no action on a selected
component or on the chosen semistable model is needed. We apply the
extension theorem componentwise and take the direct sum. We use this
semistable pair itself: an arbitrary further resolution need not
preserve reducedness.

For application of the cited theorem, $H$ consists of the horizontal
logarithmic boundary.  The central fiber is not included in $H$:
it appears in the relative logarithmic complex, with $dt/t$ in the
base quotient.  Write $f:Z\to B$ for this degeneration over the disk,
and $Z_0=f^{-1}(0)$ for its central fiber.  The top relative
logarithmic sheaf is
\begin{equation}\label{rankone:log-volume}
 \omega_{Z/B}\bigl(H+(Z_0)_{\mathrm{red}}-f^*[0]\bigr).
\end{equation}
On the semistable model the central fiber is reduced, so this is
$\omega_{Z/B}(H)$.  For a smooth open fiber of dimension $d$,
Poincar\'e duality identifies
\[
 H^d_c(Z_t\setminus H_t)^*\simeq
 H^d(Z_t\setminus H_t)(d).
\]
Accordingly the dual degree-zero Hodge quotient in
\cite[Theorem~1.1(iv)]{FujinoFujisawaVMHS} is precisely the top
$F^d$ step of ordinary $H^d$.  The upper and lower canonical
extensions coincide after unipotentization.  This proves that the
canonical extension of the top step is the direct image of
\eqref{rankone:log-volume}.  The K\"ahler form of the
one-parameter limiting theorem is also explained in
\cite[Remarks~2.16, 4.4, and~4.8]{FujinoFujisawaVMHS}.

\begin{lemma}\label{rankone:extension}
The projection of the rank-one top eigenspace to its nonzero pure
weight grade is an isomorphism on canonical extensions after
unipotentization.  Consequently it is an isomorphism of parabolic
rational lines before that base change.
\end{lemma}

\begin{proof}
Let $\overline{\mathcal V}_\chi$ be the extended mixed eigensystem
over the disk.  The finite deck projector acts on the canonical
extension and its filtrations.  By the cited extension theorem,
each
\[
 \Gr_F^p\Gr^W_j\overline{\mathcal V}_\chi
\]
is locally free and has its rank on the punctured disk.  The
filtrations on the weight quotients are induced filtrations.
Thus the exact sequences
\[
 0\longrightarrow F^dW_{j-1}\overline{\mathcal V}_\chi
 \longrightarrow F^dW_j\overline{\mathcal V}_\chi
 \longrightarrow F^d\Gr^W_j\overline{\mathcal V}_\chi
 \longrightarrow0
\]
have locally free quotients.  There is no step above $F^d$, also
on the extension, by these constant-rank assertions.  Exactly one
of the displayed quotients has rank one, namely that for $j=w$;
all the others have rank zero.  Induction on the weight filtration
therefore gives
\[
 F^dW_{w-1}\overline{\mathcal V}_\chi=0,\qquad
 F^dW_w\overline{\mathcal V}_\chi
   =F^d\overline{\mathcal V}_\chi,
\]
and identifies the latter line with
$F^d\Gr^W_w\overline{\mathcal V}_\chi$ across the origin, without
an additional zero or pole.  Descent gives the parabolic assertion.
\end{proof}

\subsubsection{The divisorial order of the Hodge line}

With the notation and prepared source above, we claim that the
parabolic order of the frame $s^{1/m}$ in $M$ is
\begin{equation}\label{rankone:order}
 \beta_D=\ord_D^M(s^{1/m})
     =\min_{g(E)=D}\frac{l_E+1-a_E}{a_E}.
\end{equation}
Equivalently, $\beta_D$ is the largest rational number $b$ such that
$t^{-b}s^{1/m}$ extends at $D=(t=0)$ in the parabolic sense.

It suffices to work over a transverse disk, suppressing the other
base coordinates.  Write $\tau=s^{1/m}$ and
$\Omega=\tau\wedge g^*dt$.  The multivalued total form
$\Omega/t$ has order $l_E-a_E$ at $E$.  A logarithmic extension
allows order $-1$, so twisting by $t^{-b}$ imposes
\begin{equation}\label{rankone:inequality}
 l_E+1-a_E-a_Eb\geq0.
\end{equation}
To verify this directly against the canonical extension, take a
finite base change $t=u^N$ and a semistable model of the root
cover.  At a central prime over the strict transform of $E$, let
$e$ be the total ramification index of the composite map to $X$,
including the cyclic-cover ramification.  Reducedness gives
$ea_E=\ord_{E'}(t)=N\ord_{E'}(u)=N$, where $E'$ is this new prime.
Orders of a total volume form shifted by one multiply by $e$.
Dividing its pullback by
$dt/du=Nu^{N-1}$ therefore gives
\[
 \ord(\tau\wedge du)
 =e(l_E+1)-N
 =e(l_E+1-a_E).
\]
The base twist subtracts $Nb=ea_Eb$.  Regularity in
\eqref{rankone:log-volume} is exactly
\eqref{rankone:inequality}.  By \Cref{rankone:extension}, testing
the logarithmic top form tests the chosen pure Hodge line as well.
The form being tested is the pullback of the original tautological
form on the full cover. The identification of open-fiber cohomology
above and the direct-image description of its canonical extension
identify it with that same Hodge section. Keeping all normalized
components ensures that a prime above every original $E$ is present;
the proper toroidal modifications retain its strict transform.

To prove sufficiency, continue to write $X$ for the smooth source
over the transverse disk and put $n=d+1$.  In SNC coordinates write
\[
 \Omega=h\prod_i x_i^{l_i}\,dx_1\wedge\cdots\wedge dx_n,
 \qquad t=h'\prod_i x_i^{a_i},
\]
where $h,h'$ are units and $a_i=0$ for horizontal components.
Put $c_i=l_i+1-(1+b)a_i$.  The vertical inequalities
\eqref{rankone:inequality} and the horizontal logarithmic bounds
give $c_i\geq0$.  For any additional divisorial valuation $v$,
write $A_X(v)$ for its log discrepancy over this smooth model,
that is, one plus the order of the Jacobian.  Logarithmic pullback
of the SNC coordinate divisor gives
$A_X(v)\geq\sum_i v(x_i)$, exactly as in
\Cref{setup:log-modification}.  Thus the shifted order of
$t^{-(1+b)}\Omega$ is
\[
 A_X(v)+\sum_i\bigl(l_i-(1+b)a_i\bigr)v(x_i)
 \ \geq\ \sum_i c_i v(x_i)\ \geq\ 0.
\]
These inequalities persist under finite covers, since shifted orders
multiply by ramification indices.  Further exceptional primes therefore
impose no stronger condition.
Horizontally, shifted orders are strictly positive outside the
designated log boundary; on the root cover these are positive
integers, so there are no poles there.  Conversely, covering primes
above strict transforms detect every
inequality in \eqref{rankone:inequality}.  Taking their minimum
proves \eqref{rankone:order}.

\begin{remark}\label{rankone:collision}
The use of all primes on the resolved model in
\eqref{rankone:order} matters even for a simple logarithmic
family.  On $\mathbf P^1_x$ over a disk, remove the horizontal
divisor $x^2=t$ and use $\tau=dx/(x^2-t)$.  After $t=u^2$ and
$x=uz$, one has $\tau=u^{-1}dz/(z^2-1)$, so the parabolic order
is $-1/2$.  Resolving the tangency of the horizontal divisor with
$t=0$ introduces a prime with $a_E=2$ and $l_E=0$, which gives
exactly $-1/2$ in \eqref{rankone:order}.  The original central
prime alone would not detect it.
\end{remark}

\subsubsection{The boundary and the comparison of sections}

Define
\begin{equation}\label{rankone:alpha-T}
 T_D=\alpha_D-\beta_D.
\end{equation}
Changing $s$ to $v s$, for a meromorphic base function $v$, adds
$a_E\ord_D(v)/m$ to $l_E$.  Hence it adds
$\ord_D(v)/m$ to both $\alpha_D$ and $\beta_D$ and leaves
$T_D$ unchanged.  This also checks the transformation rule for
the rational Hodge frame in \eqref{rankone:order}.  Passing to
a higher divisible generating degree has the same effect: on the
generic fiber its generator is a base multiple of a power of $s$.
Thus the construction is independent of these choices.

For the boundary inequality, write
\[
 x_E=\frac{l_E+\delta_E}{a_E},\qquad
 c_E=1-\frac{1-\delta_E}{a_E}.
\]
Then $0\leq c_E\leq1$, and
\[
 T_D=\min_E x_E-\min_E(x_E-c_E).
\]
For any finite collections of real numbers $x_E,c_E$, this
difference lies between $\min_Ec_E$ and $\max_Ec_E$.
Moreover the inf-multiplicity convention gives
\[
 \min_Ec_E
 =1-\max_{g(E)=D}\frac{1-\delta_E}{a_E}
 =B(g,\Delta)_D.
\]
This proves \eqref{rankone:boundary}.  If all $\delta_E=1$,
all $c_E=1$, so $T_D=1$.

Choose the discriminant to include the vertical boundary and all
loci where the smooth log construction fails.  Off it the family
and the generator line are smooth log-relative data, and there is
no vertical boundary.  There $\alpha_D=\beta_D$; in a
nonvanishing generator frame both are zero.  Thus $T$ has finite
support in the chosen discriminant, which may be resolved to be
SNC.  On subsequent modifications one pulls back the variation and
its parabolic line and recomputes $\alpha_D$ and
\eqref{rankone:alpha-T}; the order computation is unchanged on the
resulting models.

\begin{proof}[Completion of the proof of \Cref{rankone:comparison}]
Take $q$ divisible by $m$ and by the denominators of all rational
data.  In a local base frame, a meromorphic section of
$q(K_W+T+M)$ is represented
by
\[
 \varphi\,\omega^q s^{q/m}.
\]
The order of the frame in $T+M$ is
$q(T_D+\beta_D)=q\alpha_D$.  Consequently this section is
regular at $D$ exactly when \eqref{rankone:base-bound} holds.
The source-side order test already proved identifies this with
the boundary bounds at every source prime dominating $D$.
Along horizontal primes the required bound is automatic because
$s^{q/m}$ is a logarithmic
pluricanonical section on the general fiber.  We have proved the
exact section comparison in codimension one on the base, and at
all source primes except those contracted by $g$ into codimension
at least two.

Conversely, a section of $q(K_X+\Delta)$ restricts on a general
$g$-fiber to a multiple of $s^{q/m}$.  Their ratio is a
meromorphic function constant on the connected general fibers,
and therefore descends meromorphically to $W$.  The same
divisorial inequalities make it a section of $q(K_W+T+M)$
outside codimension two; smoothness of $W$ extends it uniquely
across that set.  This proves the converse comparison in all
divisible degrees under consideration.

For the forward direction, a section on $W$ gives a meromorphic
adjoint form on $X$ whose only possible remaining poles lie on
$g$-contracted prime divisors.  By hypothesis these divisors are
$\pi$-exceptional.  Push the form to $X_0$.  At every prime of
$X_0$ its transform was among the tested primes of $X$, so the
pushed form satisfies the $\Delta_0$ pole bound in codimension
one.  Since $X_0$ is smooth and the adjoint multiple is a line
bundle, it extends across codimension two.  Finally
\Cref{setup:log-modification} pulls it back to a section of
$q(K_X+\Delta)$, proving \eqref{rankone:sections}.

If the diagram is over $Y$, apply this argument over an arbitrary
open subset $U\subset Y$.  Its inverse image in $X$ is
$\pi^{-1}(f_0^{-1}U)=g^{-1}(h^{-1}U)$, and $\pi$ is still
proper over the smooth open space $f_0^{-1}U$.  The exceptional
images still have codimension at least two.  The construction and
the extension of sections commute with restriction to $U$,
giving \eqref{rankone:relative-sections}.  Relative canonical
versions, when the auxiliary base is smooth, follow by the
projection formula.  All comparisons multiply by the same relative
generator, with its powers in the common divisible grading, so they
preserve products and ratios of sections. The projection formula also
gives the identity after any pulled-back base twist, in degrees clearing
the denominator for a rational twist. This completes
\Cref{rankone:comparison}.
\end{proof}

\subsection{Retaining the original base multiplicities}

\begin{lemma}[Composition of inf-multiplicity bounds]
\label{application:composition}
Suppose that \(X\xrightarrow{g}W\xrightarrow{h}Y\) is a diagram of
surjective maps between smooth spaces, and \(\Delta,T,C\) are rational
boundaries on these spaces.  If
\[
 T\geq B(g,\Delta),\qquad C\leq B(h\circ g,\Delta),
\]
then \(B(h,T)\geq C\).  If, in addition, every prime divisor of \(X\)
mapped into codimension at least two in \(Y\) has coefficient one in
\(\Delta\), and \(T\leq1\), then every prime divisor of \(W\) mapped
into codimension at least two in \(Y\) has coefficient one in \(T\).
\end{lemma}

This is the inf-multiplicity composition comparison of
\cite[Proposition~3.13(1)]{CampanaOrbifolds}; we include its order
calculation to track coefficient-one exceptional divisors.

\begin{proof}
Let \(D\subset Y\) and \(V\subset W\) be prime divisors with
\(h(V)=D\).  Write \(b_V=\ord_V(h^*D)\).  If \(E\subset X\) is a
prime divisor dominating \(V\), put \(a_E=\ord_E(g^*V)\).  At their
generic points orders multiply, so
\[
 \ord_E((h\circ g)^*D)=b_Va_E.
\]
The hypothesis on \(C\) implies
\(b_Va_E m_\Delta(E)\geq m_C(D)\) for every such \(E\).  Taking the
minimum over them and then using the hypothesis on \(T\) gives
\[
 b_Vm_T(V)\ \geq\ b_Vm(g,\Delta;V)\ \geq\ m_C(D).
\]
Taking the minimum over \(V\) proves the first assertion.  All these
inequalities are valid with the prescribed convention for infinity.

For the second assertion, let \(V\) be mapped into codimension at least
two in \(Y\).  Every prime \(E\) dominating it has
\(\Delta_E=1\), so \(m(g,\Delta;V)=\infty\).  It follows that
\(B(g,\Delta)_V=1\), and hence \(T_V=1\).
\end{proof}
